% ============================================================================
% Chapter 38: Email and Online Communication
% ============================================================================
\chapter{Email and Online Communication}
\label{ch:email}

\begin{objectives}
  \item Create and send emails with proper etiquette
  \item Attach files to an email
  \item Understand CC, BCC, and Reply All
  \item Recognise phishing emails and online scams
\end{objectives}

\bytetip[tip]{Email is one of the most important communication tools in the world. Learning to write proper emails now will help you throughout school, college, and your career!}

\section{What Is Email?}
\label{sec:what-email}
\index{email}

\hl{Email} (Electronic Mail) is a way to send and receive digital messages over the Internet. Popular services include \textbf{Gmail}, \textbf{Outlook}, and \textbf{Yahoo Mail}.

\section{Parts of an Email}
\label{sec:email-parts}

\begin{theorybox}[Parts of an Email]
\begin{tabular}{ll}
  \textbf{To} & The recipient's email address \\
  \textbf{CC} & Carbon Copy --- send a copy to others (visible to all) \\
  \textbf{BCC} & Blind Carbon Copy --- copy to others (hidden from all) \\
  \textbf{Subject} & A short title describing the email's purpose \\
  \textbf{Body} & The main message content \\
  \textbf{Attachments} & Files sent along with the email \\
  \textbf{Signature} & Your name and contact info at the bottom \\
\end{tabular}
\end{theorybox}

\section{Email Etiquette}
\label{sec:etiquette}
\index{email!etiquette}

\begin{enumerate}[label=\textcolor{ModuleBlue}{\textbf{\arabic*.}}, itemsep=3pt]
  \item Always write a clear, specific \textbf{subject line}
  \item Start with a \textbf{greeting} (Dear Sir/Ma'am, Hello, Hi)
  \item Keep the message \textbf{short and clear}
  \item Use \textbf{polite language} --- no ALL CAPS (it looks like shouting!)
  \item Always \textbf{proofread} before sending
  \item End with a \textbf{sign-off} (Kind regards, Thank you, Best wishes)
\end{enumerate}

\begin{warnbox}
\textbf{Phishing alert!} Watch out for emails that:
\begin{itemize}[leftmargin=*, itemsep=2pt]
  \item Ask for your \textbf{password} or personal information
  \item Contain \textbf{urgent threats} (``Your account will be deleted!'')
  \item Come from \textbf{unknown senders} with suspicious links
  \item Have \textbf{spelling mistakes} in the sender's address (e.g., ``gooogle.com'')
\end{itemize}
When in doubt, \textbf{don't click any links} --- ask a parent or teacher first!
\end{warnbox}

\bytetip[warn]{Once you send an email, you \textbf{can't un-send it} (well, Gmail gives you about 30 seconds). Always double-check the recipient and re-read your message before hitting Send!}

\begin{funfact}
The first email was sent in 1971 by \textbf{Ray Tomlinson}. He chose the \textbf{@} symbol to separate the user name from the computer name. He later admitted he couldn't remember what the first email said --- probably ``QWERTYUIOP'' or something similar as a test!
\end{funfact}

\begin{reviewqs}
  \item What does CC stand for, and how is it different from BCC?
  \item Name three things you should include in a formal email.
  \item What are three signs that an email might be a phishing scam?
  \item Why should you always write a clear subject line?
  \item How do you attach a file to an email?
\end{reviewqs}

\begin{challenge}
Write a \textbf{formal email} to your teacher asking for extra time on a homework assignment. Include: proper greeting, clear subject line, polite body with a reason, and a professional sign-off. Save it as a draft or send it for real!
\end{challenge}
